\scanpage{58}
\begin{center}4 - 35 -\hfill n\textsuperscript{o} 392\end{center}

libre de $R$ tel que $R/M$ soit un $B$-module de torsion ; comme $M_{[A]}$ est somme directe de $r_B(R)$ $A$-modules isomorphes à $B$, on a, (prop. 16, (i)), $c_A(M_{[A]})= {}=r_B(R).c_A(B)$. D'autre part, $c_A((R/M)_{[A]})=-c_A(N(\unsure{\chi|\gamma}_B(R/M)))$ en vertu de (19) ; mais par définition de l'homomorphisme $N:C(B)\longrightarrow C(A)$, on a $c_A(N(d))=N(c_B(d))$ pour tout $d\in D(B)$, et comme $c_B(\unsure{\chi|\gamma}_B(R/M))=-c_B(R)$ par définition, on a finalement $c_A((R/M)_{[A]})=N(c_B(R))$ ; il suffit alors d'appliquer la prop. 16, (i) pour obtenir (20).

PROPOSITION 19.- \uline{Soit} $R$ \uline{un} $B$-\uline{module de type fini}. \uline{Pour que} $R$ \uline{soit réflexif}, \uline{il faut et il suffit que} $R_{[A]}$ \uline{soit un} $A$-\uline{module réflexif}.

Il est clair que pour que $R$ soit un $B$-module sans torsion, il faut et il suffit que $R_{[A]}$ soit un $A$-module sans torsion. On peut donc supposer que $R$ est un réseau de $W=R\otimes_B L$ par rapport à $B$. Nous utiliserons le lemme suivant :

LEMME 4.- \uline{Soit} $W$ \uline{un espace vectoriel de dimension finie sur} $L$ \uline{et soit} $R$ \uline{un réseau de} $W$ \uline{par rapport à} $B$. \uline{Pour tout} $\underline{p}\in P(A)$, \uline{on a alors} $(R_{[A]})_{\underline{p}}=\bigcap_{\underline{P}\mid\underline{p}}R_{\underline{P}}$.

En effet, si $S=A-\underline{p}$, les idéaux premiers de l'anneau $S^{-1}B$ sont engendrés par les idéaux premiers de $B$ ne rencontrant pas $S$, autrement dit les idéaux $\underline{P}_i$, $(1\le i\le m)$, au-dessus de $\underline{p}$ et l'idéal $(0)$ ; cela montre que $S^{-1}B$ est un anneau semi-local dont les idéaux maximaux sont les $\underline{m}_i=\underline{P}_i(S^{-1}B)$ pour $1\le i\le m$ ; en outre l'anneau local $(S^{-1}B)_{\underline{m}_i}$ est isomorphe à $B_{\underline{P}_i}$, (chap. II, § 2, n\textsuperscript{o} 5, prop. 11), donc est un anneau de valuation discrète. L'anneau $S^{-1}B$ est par suite un anneau de Dedekind, (§ 3, n\textsuperscript{o} 2, th. 1, e)), et comme il est semi-local, il est principal, (§ 3, n\textsuperscript{o} 2, prop. 1). Cela étant, $(R_{[A]})_{\underline{p}}$ est égal à $S^{-1}R$, considéré comme $A_{\underline{p}}$-module ; d'après ce qui précède, $S^{-1}R$ est un réseau \uline{libre} de $W$ par rapport à $S^{-1}B$ et on peut par suite lui appliquer le th. 2 du n\textsuperscript{o} 2, qui donne $S^{-1}R=\bigcap_i(S^{-1}R)_{\underline{m}_i}$ ; mais $(S^{-1}R)_{\underline{m}_i}=R_{\underline{P}_i}$, ce qui prouve le lemme.

\margin{57}
